% Owner A: Ziyan Song.
\section{Rationale and Research Basis}
TFEB is a member of the MiT/TFE transcription-factor family and a major transcriptional regulator of the autophagy--lysosome pathway. Under nutrient-rich conditions, mTORC1-dependent phosphorylation favors cytoplasmic retention of TFEB. During nutrient starvation or mTOR inhibition, TFEB accumulates in the nucleus and activates genes involved in lysosomal biogenesis and cellular clearance \cite{Settembre2012,RoczniakFerguson2012}. Nutrient-dependent control of nuclear export also contributes to TFEB distribution \cite{Napolitano2018}.

This project connects molecular cloning with a visible cell-biological phenotype. Its primary readout is the nuclear-to-cytoplasmic distribution of a fluorescent TFEB fusion. The study will test whether the sequence-verified reporter exhibits the expected nutrient-dependent redistribution in HeLa cells.

\subsection{Research Objectives and Hypothesis}
The first objective is to construct and sequence-verify an EGFP-tagged TFEB mammalian expression plasmid. The second is to compare its subcellular localization under normal culture, starvation, and mTOR-inhibitor conditions. The hypothesis is that starvation and mTOR inhibition increase the nuclear-to-cytoplasmic fluorescence ratio relative to the vehicle-treated control.
